\chapter{Foundational Theorem Dossier}
\label{app:perverse-foundations}

This appendix records every theorem whose full SGA-scale construction is not
reproduced line by line.  Each entry states the exact form used, its required
compatibility, and the point at which it enters.  None contains the
Deligne--Laumon weight theorem or smooth-proper purity.

\section{Scheme theory}

\begin{longtable}{
 >{\raggedright\arraybackslash}p{0.22\textwidth}
 >{\raggedright\arraybackslash}p{0.38\textwidth}
 >{\raggedright\arraybackslash}p{0.30\textwidth}}
\toprule
Input & Exact statement used & Role\\
\midrule
\endhead
Nagata compactification &
A separated finite-type morphism of noetherian schemes factors as an open
immersion followed by a proper morphism; two factorizations admit a common
refinement. &
Construction and independence of \(Rf_!\).\\

Relative Noether normalization &
After shrinking integral source and target, a dominant affine finite-type
morphism factors as a finite map to affine space over the target. &
Reduction of global weights to relative affine lines.\\

Stratification &
A constructible complex is lisse on a finite locally closed stratification;
over the perfect field \(\Fq\), every reduced finite-type scheme has a dense
smooth open. &
Perverse gluing and Noetherian induction.\\

Curves &
A smooth curve has a smooth proper completion with finite boundary; after a
finite constant-field extension it admits a finite map to \(\PP^1\) with
controlled boundary. &
Placement of the curve problem on \(\A^1\).\\
\bottomrule
\end{longtable}

\section{Constructible adic derived categories}

For every finite-type \(X_0/\Fq\) and \(\ell\ne\operatorname{char}\Fq\), the
proof uses a bounded constructible adic category
\(\Dbc(X_0,\Qbarl)\) with stalks, costalks, tensor product, internal Hom, and
the six operations.  It is obtained from compatible
\(\mathbf Z/\ell^n\)-systems by quotienting negligible systems and inverting
\(\ell\).  The required interfaces are:

\begin{enumerate}
\item functorial \(Rf_!,Rf_*,f^*,Rf^!,\otimes^L,R\mathcal Hom\), with coherent
      composition isomorphisms;
\item localization triangles for an open--closed decomposition;
\item projection formula and compact-support Kunneth;
\item constructibility and a finite cohomological-amplitude bound;
\item compactification independence for \(Rf_!\).
\end{enumerate}

All isomorphisms are required in the Weil category, so they commute with
geometric Frobenius.  Replacing a canonical isomorphism with definitional
equality would make compactification choices leak into later statements.

\section{Base change, purity, and duality}

\begin{longtable}{
 >{\raggedright\arraybackslash}p{0.21\textwidth}
 >{\raggedright\arraybackslash}p{0.43\textwidth}
 >{\raggedright\arraybackslash}p{0.26\textwidth}}
\toprule
Theorem & Form used & Compatibility\\
\midrule
\endhead
Proper base change &
For proper \(f:X\to S\),
\[
 (Rf_*\cK)_{\bar s}\simeq
 R\Gamma(X_{\bar s},\cK|_{X_{\bar s}}).
\]
&
Base automorphisms and Frobenius.\\

Proper-support base change &
For separated finite-type \(f\), the corresponding identity with \(Rf_!\)
and compactly supported fiber cohomology. &
Composition and localization.\\

Smooth purity &
For smooth \(f\) of pure relative dimension \(d\),
\[
 f^!\Ql\simeq\Ql(d)[2d].
\]
&
Products, smooth base change, and the affine-space fundamental class.\\

Verdier duality &
\[
\begin{gathered}
 \D_YRf_!\simeq Rf_*\D_X,\\
 R\mathcal Hom(Rf_!\cK,\cL)\\
 \simeq
 Rf_*R\mathcal Hom(\cK,f^!\cL).
\end{gathered}
\]
&
Adjunction, projection formula, compactification.\\

Poincare duality &
For smooth \(d\)-dimensional \(X\),
\[
 H_c^i(X,\cF)\times H^{2d-i}(X,\cF^\vee(d))
 \to\Qbarl
\]
is perfect. &
Galois equivariance; untwisted paired eigenvalues multiply to \(q^d\).\\
\bottomrule
\end{longtable}

Absolute cohomological purity for a regular immersion is used in the local
costalk computation proving that \(\cL[d]\) is perverse.  It is not the
arithmetic purity of Frobenius eigenvalues.

\section{Trace and the Artin--Schreier kernel}

\begin{theorem}[Grothendieck--Lefschetz trace formula: imported interface]
\label{thm:dossier-trace}
For \(X_0/\Fq\), a constructible Weil complex \(\cK_0\), and \(n\geq1\),
\[
 \sum_{x\in X_0(\mathbf F_{q^n})}
 \Tr(F_x\mid\cK_{\bar x})
 =
 \sum_a(-1)^a
 \Tr(F^n\mid H_c^a(X,\cK)).
\]
For a closed-point formulation, a point of degree \(e\mid n\) contributes
\(F_x^{\,n/e}\), not \(F_x^n\).  The formula is additive in triangles and
compatible with proper pushforward and extension by zero.
\end{theorem}

For a nontrivial additive character
\(\psi:\mathbf F_p\to\Qbarl^\times\), the Lang isogeny
\(t\mapsto t^p-t\) supplies the rank-one Artin--Schreier sheaf
\(\cL_\psi\) on \(\A^1\).  Its trace function is
\(x\mapsto\psi(\Tr_{\mathbf F_{q^n}/\mathbf F_p}x)\), it is lisse and pure of
weight \(0\), and character orthogonality gives the delta-kernel calculation
used in Fourier inversion.

\section{Perverse foundations}

The perverse layer consumes the following non-arithmetic statements from BBD
Sections 1--4:

\begin{enumerate}
\item gluing of bounded \(t\)-structures along an open--closed recollement;
\item the middle support/cosupport \(t\)-structure on \(\Dbc(X)\);
\item Verdier self-duality;
\item affine Artin vanishing and affine/quasi-finite perverse amplitudes;
\item finite length of \(\Perv(X)\) and classification of its simple objects.
\end{enumerate}

Their constructions are given in
\cref{chap:perverse-t-structure,chap:intermediate-extension}.  No theorem from
BBD Sections 5--6 concerning weights, semisimplicity, or decomposition is an
input to the Fourier/curve proof.

\section{Geometric reductions used by Fourier}

The curve proof additionally uses:

\begin{enumerate}
\item finite maps of curves and proper base change;
\item Grothendieck--Ogg--Shafarevich to control Euler characteristic and
      boundary conductors;
\item local monodromy filtration for a pure lisse sheaf on a punctured trait;
\item the rank-one purity theorem and the real-sheaf criterion, established
      upstream of the general direct-image theorem.
\end{enumerate}

The local monodromy statement is: if \(\cL\) is pure of weight \(w\), then
\(\operatorname{Gr}^M_r\) of the monodromy filtration is pure of weight
\(w+r\).  It is a local theorem on a trait; using global smooth-proper purity
to prove it would be circular, so the notes use Deligne's local
monodromy argument.

\section{What is proved rather than imported}

For clarity, the following are outputs of the book:

\begin{itemize}
\item the real-sheaf mixedness theorem in the form needed by Fourier;
\item Fourier inversion, Plancherel, and perverse \(t\)-exactness;
\item the curve weight estimate;
\item the general \(Rf_!\) upper-weight theorem;
\item proper direct-image purity and purity of intersection cohomology;
\item geometric semisimplicity and decomposition;
\item all clauses of the classical Weil conjectures, including integrality
      and \(\ell\)-independence.
\end{itemize}

This ledger is the boundary between a one-volume proof and the SGA-scale
construction of its cohomological backend.
